\section{Introduction}
\label{sec:introduction}

Urban agglomerations in the Global South represent the ground zero of contemporary climate vulnerability. Rapid, often uncoordinated demographic expansion and extensive land-use conversion have created dense metropolitan fabrics where environmental hazards intersect with high infrastructural and population exposure \citep{ipcc2022ar6}. In these vulnerable urban centers, environmental stress does not manifest as isolated, single-hazard phenomena. Rather, extreme heat, hydrological deficits, toxic atmospheric particulate matter ($\text{PM}_{2.5}$), and vegetative canopy loss routinely converge across space and time, producing compound multi-hazard events whose collective societal and physiological impacts far exceed the sum of their individual parts \citep{zscheischler2018future}.

\subsection{The Fragmented Landscape of Existing Urban Risk Indices}
Despite urgent calls by the Intergovernmental Panel on Climate Change (IPCC) to address compound and cascading risks, current academic research and municipal decision-making frameworks remain deeply fractured into single-domain silos:

\begin{enumerate}
    \item \textbf{The Purely Biophysical / Meteorological Silo:} Atmospheric and climatological frameworks focus predominantly on physical weather extremes---such as the Universal Thermal Climate Index (UTCI), the Standardized Precipitation-Evapotranspiration Index (SPEI), or municipal Air Quality Indices (AQI). While these indices capture atmospheric dynamics, they treat the urban environment as an isotropic, featureless surface. They fail to account for micro-scale spatial heterogeneity in building morphology, impervious surface fractions, and localized exposure, which dictate how atmospheric stress is experienced on the ground.
    \item \textbf{The Purely Urban / Morphological Silo:} Conversely, urban planning and geographical vulnerability assessments---such as the Social Vulnerability Index (SoVI) or structural built-density mappings---focus almost exclusively on stationary infrastructure and census demographics. These frameworks capture spatial form but treat climate and environmental conditions as temporally static or spatially uniform background constants. Consequently, they completely miss high-frequency seasonal climate cycles, acute heatwaves, and seasonal atmospheric stagnation events.
    \item \textbf{The Linearity and Substitutability Fallacy:} When environmental and urban metrics are combined into composite indicators, existing models almost universally enforce naive linear additive aggregation ($\sum w_i x_i$). Under international composite indicator guidelines \citep{oecd2008handbook}, linear aggregation imposes an assumption of perfect substitutability: deficits in one critical domain (e.g., lethal particulate inhalation or extreme water scarcity) are mathematically offset by favorable scores in another (e.g., proximity to suburban green space). In biological, ecological, and civil infrastructure systems, this assumption is fundamentally invalid: concurrent extreme hazards trigger non-linear, threshold-activated systemic failures.
\end{enumerate}

\subsection{Methodological Distortions in Composite Environmental Index Design}
Beyond structural silos, existing composite environmental indices suffer from three prevalent mathematical distortions:
\begin{itemize}
    \item \textit{The Seasonal Normalization Trap:} Panel indices frequently standardize variables by calendar month using dynamic sample $z$-scores ($z = (x - \mu_{\text{month}})/\sigma_{\text{month}}$). This mathematically forces the monthly city-wide average of every hazard to an artificial mid-scale score of $0.50$, effectively subtracting the primary phenomenon of interest: true annual climate seasonality.
    \item \textit{The Boundary Ceiling Distortion:} Naive non-linear coupling models that add an interaction penalty directly to a linear base ($\min(1.0, \text{Linear} + \text{Compound})$) suffer from severe boundary truncation. In multi-city panels, this hard ceiling clamps a large fraction of observations at exactly $1.0000$, destroying relative spatial rank discrimination among the most severely impacted neighborhoods.
    \item \textit{Dimensionality Collapse in PCA Weighting:} While Principal Component Analysis (PCA) is frequently employed for objective weight derivation, environmental hazard matrices in diverse biomes rarely exhibit high multi-collinearity. When the leading component (PC1) accounts for less than $40\%$ of overall variance, PCA loading vectors collapse the weights of critical, non-collinear hazards (such as atmospheric aerosols and stationary built infrastructure) down to near zero, biasing the composite index toward acting merely as a regional weather proxy.
\end{itemize}

\subsection{Contributions of this Work}
To resolve these theoretical and methodological challenges, this paper introduces the \textbf{Compound Sustainability Stress Index (SSI)}, a unified spatio-temporal framework that bridges satellite earth observations, atmospheric reanalysis, and micro-scale urban built environment vectors into a continuous, non-linear multi-hazard diagnostic. 

Grounded in an empirical longitudinal panel of \textbf{334,332 discrete hexagonal micro-units} (Uber H3 Resolution~8, $\sim 0.74\text{ km}^2$) monitored continuously across 36 months (2020--2022) across seven contrasting Indian metropolitan biomes, this paper establishes three primary contributions:
\begin{enumerate}
    \item \textbf{Cross-Domain Integration on an Invariant Spatial Grid:} We bridge five native physical drivers spanning meteorological heat, hydrological water deficit, photosynthetic vegetation loss, atmospheric $\text{PM}_{2.5}$, and built-environment physical exposure onto an objective Discrete Global Grid System, eliminating administrative boundary distortions (MAUP).
    \item \textbf{Smooth Asymptotic Non-Linear Compound Coupling:} We formulate a bivariate joint-tail compound hazard amplifier with dynamic damping that models extreme multi-hazard synergy ($+25.3\%$ stress escalation during concurrent tail events) while guaranteeing that scores remain strictly bounded within $[0, 1]$ with zero ceiling clipping artifacts.
    \item \textbf{Pan-India Typology and Pre-Registered Ground-Truthing:} We establish an eight-part validation battery that verifies indicator orthogonality, spatial coherence, cross-biome differentiation ($F = 6,888$), convergent construct validity against independent satellite thermal observations ($\rho = 0.72\text{--}0.88$), and precise alignment with recorded historical climate disasters. Unsupervised clustering delineates six operational stress archetypes to guide municipal adaptation policy.
\end{enumerate}
